\section{Lesson 5: What If I Say No?}

Once you have stated what a metaphysical passage claims, fixed its terms, sorted the warrant behind each premise, located the burden-bearing sentence, and bounded what follows, one question turns passive reading into genuine philosophical judgment: \textbf{Denial}. What happens to a reader who says no to the burden-bearing premise?

Asking what denial costs converts a wall of authoritative assertion into a concrete choice with an inspectable price. It also protects Catholic philosophy from a bad defence. Defenders of scholastic arguments often declare that denying a substantive premise is ``absurd'' or ``incoherent.'' Sometimes that charge holds, and exhibiting a contradiction is decisive. Far more often, rejection commits the denier to a competing metaphysical account, leaves an explanatory debt, or simply declines a technical vocabulary. Claiming contradiction where none exists misleads sympathetic readers, blocks honest engagement with rivals, and erodes credibility where contradiction truly occurs.

\subsection{The Five Kinds of Denial}

Rejecting a burden-bearing premise produces one of five distinct outcomes. Each takes shape on real denials of real metaphysical premises.

\subsubsection*{Contradiction}

A denial incurs \textbf{contradiction} when rejecting the premise forces the denier to affirm and deny the same proposition in the same respect, or to retract something already granted. This is the strongest verdict against a denier and the rarest in metaphysics. Claiming it requires \emph{exhibiting} the contradiction step by step, never merely alleging that incoherence lurks nearby.

Take Aquinas's premise in the Second Way (\emph{Summa Theologiae} I, q.\,2, a.\,3, corp.): nothing can be the efficient cause of itself, ``for so it would be prior to itself, which is impossible.'' Grant that an efficient cause produces something into being and must exist in order to act. Anyone who then asserts that a thing efficiently caused its own existence must hold that it already existed in order to act and did not yet exist in order to be produced---existing and not existing in the same respect. Aristotle's elenctic defence of non-contradiction (\emph{Metaphysics} Gamma 3--4, 1005b19--20, 1006a28--1009a5) stands in the same rare class: signifying anything definite already presupposes that signifying it is not identical with not signifying it.

\subsubsection*{Conflict with experience}

A denial incurs \textbf{conflict with experience} when rejecting the premise contradicts what we plainly and commonly encounter, such as the First Way's opening premise (\emph{ST} I, q.\,2, a.\,3, corp.): ``It is certain, and evident to our senses, that in the world some things are in motion.'' Denying that change occurs clashes directly with everyday perception. Even here, however, denial is weaker than contradiction. Parmenides in antiquity and radical four-dimensionalists today reinterpret apparent change as static slices of an unchanging reality. Such reinterpretations bear a steep empirical burden and deserve sharp scrutiny, yet they must be examined rather than ridiculed.

\subsubsection*{Explanatory cost}

A denial incurs an \textbf{explanatory cost} when rejecting the premise leaves unexplained a real phenomenon or regularity that the premise accounted for. Deny the substantive reading of \emph{nemo dat quod non habet} (``nothing gives what it does not have'')---accepting that wholly novel features emerge from causes with no foundation for them---and you incur neither contradiction nor conflict with bare observation. What you incur is an unanswered question about why effects regularly correspond to their causes.

Hold one boundary firm: \textbf{an explanatory cost is a real disadvantage, never a refutation.} A denier may openly accept a brute stopping point, dispute whether the phenomenon requires that kind of explanation, or offer a different explanation. Explanatory cost differs sharply from a rival framework: an explanatory cost is a specific debt left on the table, whereas a rival framework is a developed competing ontology that reorganizes the whole domain, usually trading one set of explanatory costs for another.

\subsubsection*{Rival framework}

A denial commits the denier to a \textbf{rival framework} when rejection belongs to a developed competing metaphysical account of the same subject matter---for instance, replacing Aristotelian potency with microstructural dispositions, or replacing Aquinas's real distinction of essence and existence with a conceptual distinction. The dispute is now between two coherent positions compared on their principles and trade-offs, not between truth and logical error. State this plainly from the start: \textbf{discovering a coherent rival framework is not a defeat of the tradition.} Open questions between developed accounts are the normal condition of serious metaphysics.

\subsubsection*{Non-assent}

A denial is \textbf{non-assent} when the denier merely declines a vocabulary, conceptual scheme, or project, incurring no contradiction, no conflict with plain experience, and no explanatory loss. He has made no mistake; he has simply declined to join a way of describing things. When a chemist declines to describe freezing water as losing the accidental form of liquidity while keeping its substantial form---preferring to speak strictly of kinetic energy and hydrogen-bond geometry without erecting an anti-hylomorphic metaphysics---he stands in \textbf{non-assent}.

\subsection{Where Traditional Defences Most Often Overreach}

The central discrimination of this method is separating \textbf{rival framework} and \textbf{non-assent} from \textbf{contradiction}. Test for it by pushing the hostile-reader question one step further: construct the denier explicitly, give him a voice, and state how he handles the everyday phenomena. If his position stands without affirming and denying the same thing in the same respect, denial is not contradiction.

Watch how easily rhetoric overreaches in a familiar scholastic argument for bodily unity:

\begin{quote}
Every natural body is composed of substantial form and prime matter, because a genuine substance is one unified thing rather than a heap of parts, and what makes it one unified substance is its substantial form. Whoever denies hylomorphism therefore contradicts himself: he admits that the oak tree before his eyes is one single tree while denying the very principle that makes it one.
\end{quote}

Run the explicit-denier test. A biologist who declines the vocabulary of ``substantial form'' and ``prime matter'' agrees completely that the oak is one living organism rather than a pile of sawdust; he describes that unity through integrated cellular metabolism, genetic regulation, and homeostasis. He never asserts both that the oak is one organism and that it is not. What the passage's rhetoric brands as self-contradiction is in fact \textbf{non-assent} to an Aristotelian vocabulary (or, if developed into an ontology of organisms, a \textbf{rival framework}).

Say this without hedging: \textbf{treating non-assent to a vocabulary as self-contradiction is where traditional defences most often overreach and damage their own credibility.} It misleads trusting readers into mistaking open questions for settled logic, keeps real rivals off the table, and makes sharp critics distrust the tradition even when contradiction genuinely holds.

\subsection{Counterexample Discipline, Rivals, and Honest Replies}

Fairness binds critics as tightly as defenders. A counterexample touches a claim only when it respects every \emph{per se} restriction, every ``in this respect'' clause, and every fixed technical sense. Dropping a qualification misses the claim.

Consider the favourite classroom counterexample to Aquinas's First Way (\emph{ST} I, q.\,2, a.\,3, corp.): ``whatever is in motion is put in motion by another.'' A critic throws a baseball across a field; once it leaves the pitcher's hand, it keeps flying with no hand pushing it, and in empty space it would coast forever. Yet Aquinas explicitly defines motion (\emph{motus}) two sentences later as ``the reduction of something from potentiality to actuality,'' illustrated by fire heating wood. Before the baseball touches Aquinas's premise, the critic must ask whether uniform inertial coasting in empty space counts as a continuous transition from potency to act, or whether uniform velocity is a steady state rather than a potency-to-act change. Only after settling that distinction does a genuine objection emerge.

To test a premise honestly, we adopt the tradition's own form of objection and reply---the discipline of medieval \emph{disputatio}, where the strongest objections stand openly beside the reply. (Our fivefold classification of denial remains a modern reader's diagnostic for private study, not Aquinas's own taxonomy.) Three habits keep that form honest:

\begin{itemize}
\item \textbf{State every rival at its strongest}, in the form its ablest defender would accept, and answer it on those exact terms without weakening it in the reply.
\item \textbf{Expose and refuse the pseudo-reply.} A pseudo-reply restates an opponent's objection inside house vocabulary and declares it refuted without engaging its substance---for example, answering a microstructural rival by saying ``to deny real potency is to deny change itself, which is absurd.'' That merely equates the opponent's refusal of Aristotelian potency with denying observable change.
\item \textbf{Honour concession and open objections.} Where an objection genuinely succeeds in part, grant what it wins and narrow the conclusion accordingly. Where no adequate reply is available, leave the objection honestly standing rather than offering a bad reply.
\end{itemize}

\subsection{Worked Passage 1: Denying ``Whatever Is Moved Is Moved by Another''}

Return to the First Way of St.\ Thomas Aquinas, \emph{Summa Theologiae} I, q.\,2, a.\,3, corp., whose full text, meaning, warrant, burden, and bounded consequence we established in Lessons 3 and 4. Recall the exact burden-bearing premise isolated there (alongside the instrumental-series regress claim): \emph{that the potency invoked by change is a real metaphysical principle inside the changing thing demanding an actual external reducer, not merely a description of possible future states or physical structure}. Recall also our bounded consequence---at least one first mover, unmoved in the respect at issue, at the head of an essentially ordered series---and what remains \textbf{not yet established}: uniqueness, immateriality, simplicity, intelligence, goodness, or God (attributes argued separately in \emph{ST} I, qq.\,3--11 and q.\,14; \emph{SCG} I, cc.\,14--28).

\subsubsection*{Denial}

Rejecting the substantive reading of potency incurs neither \textbf{contradiction} nor bare \textbf{conflict with experience}: fire still heats wood, and nothing self-contradictory follows from rejecting an inner Aristotelian principle of potency. The accurate classification is \textbf{rival framework}, paired with mutual \textbf{explanatory costs}. State the modern rival at full strength:

\begin{quote}
I agree that fire heats wood and that nothing is hot and not-hot in the same respect at once. What I deny is that change requires a distinct metaphysical principle of ``potency'' inside wood. On a \emph{microstructural and dispositional account}, saying dry wood is ``potentially hot'' or ``flammable'' reports its actual present chemical structure together with modal and counterfactual truths about how cellulose bonds react under sufficient kinetic energy and oxygen. Everything inside the wood is fully actual, change is the lawful redistribution of actual energy, and uniform inertial motion is a conserved state requiring no ongoing mover. Or, on a sparser \emph{Humean account}, talk of capacities simply summarizes regularities across the mosaic of actual events in spacetime, with no hidden inner potencies at all.
\end{quote}

Contrast an honest response with a classic \textbf{pseudo-reply}: ``To deny real potency is to deny change itself, which is absurd; if wood had no potency to become hot, heat would spring out of pure nothingness.'' Name the failure cleanly: that pseudo-reply translates the opponent's view into house vocabulary (equating ``no Aristotelian potency'' with ``nothingness'') and cries contradiction without ever touching how actual chemical bonds and thermal energy interact.

An \textbf{honest reply and partial concession} does something far stronger. First, concede candidly that neither bare logic nor bare observation of fire refutes the microstructural or Humean accounts, and that uniform inertial motion cannot be dismissed without examining which physical transitions genuinely move from potency to act. Then press the real \textbf{explanatory cost} on the rival's own terms: the Humean must accept as brute why stable regularities hold across time without real tendencies in things, while the microstructuralist faces the question whether grounding flammability in deeper actual structures merely pushes unrealized capacities down to quantum fields and laws of nature. In return, the rival presses an explanatory cost on the Thomist: why postulate a distinct principle of potency alongside actual physical structure and conserved quantities? Both sides carry genuine philosophical debts inside an open contest of frameworks.

\subsection{Worked Passage 2: Denying the Real Distinction of Essence and Existence}

Our second passage returns to St.\ Thomas Aquinas, \emph{De ente et essentia} (\emph{On Being and Essence}), Chapter 4 (Baur/Koch Latin text, examined in full in Lesson 4), where Aquinas moves from conceiving what a man or a phoenix is without knowing whether it exists to concluding that in every reality other than at most one first reality whose essence is its existence, \emph{esse} (existence) is really distinct from essence and therefore received efficiently from another. Recall the primary \textbf{burden} fixed in Lesson 4: \emph{that in an existing finite creature, essence and existence are two really distinct metaphysical principles composing the thing, rather than two ways our minds grasp one undivided contingent reality} (alongside the ban on an endless chain of causes of existence). Recall also what remains \textbf{not yet established}: that this first cause is personal, intelligent, free, good, or provident, or that it created freely rather than necessarily.

\subsubsection*{Denial}

Denying the real distinction of essence and existence in creatures incurs neither \textbf{contradiction} nor \textbf{conflict with experience}, and it does \emph{not} imply atheism. Historically, committed Catholic scholastic thinkers rejected Aquinas's real distinction while staunchly defending creation from nothing. The honest classification is \textbf{rival framework}. State that rival scholastic position at full strength:

\begin{quote}
I agree that I can conceive a phoenix without knowing whether it exists, and I affirm wholeheartedly that every finite thing depends on God as Creator. What I deny is that an existing creature---this horse---is composed of two really distinct metaphysical principles, its essence and its existence. Outside the mind, an un-actualized horse-essence is nothing at all, and an existing essence simply is the concrete horse. Essence and existence in an existing creature therefore differ by a \emph{conceptual distinction with a foundation in reality}: our minds can consider what the horse is apart from its actual existence, and that abstraction is grounded in the real fact that the horse is contingent and depends on God's creative will rather than existing of itself.
\end{quote}

What changes under this denial is as clear as what survives. The denier incurs no contradiction and keeps divine creation intact. What disappears is Aquinas's \emph{De ente} Chapter 4 springboard: without real internal composition of essence and existence, you cannot immediately derive efficient dependence from an essence ``receiving'' an alien act of being, so dependence on God must be argued by other routes (such as contingency or finitude). Both frameworks carry reciprocal \textbf{explanatory costs}: the Thomist asks how a creature's being is genuinely limited to a finite kind without a distinct receiving potency, while the rival asks how an essence distinct from existence can have enough reality to ``receive'' and ``limit'' existence before existing.

\subsection{Worked Passage 3: Non-Assent and Rival Accounts of Perfection}

Our third passage comes from St.\ Thomas Aquinas, \emph{Summa Theologiae} I, q.\,4, a.\,1, corp.\ (``Whether God Is Perfect?'', 1920 English Dominican translation). Objection 1 notes that \emph{perfectum} literally means ``completely made,'' and God is not made; the \emph{sed contra} cites Matthew 5:48 (``Be you perfect as also your heavenly Father is perfect''), a premise resting strictly on \textbf{revelation}. In the body, Aquinas explains that early thinkers who conceived the first principle as material denied it perfection because matter is purely potential, and then gives his core argument:

\begin{quote}
Now God is the first principle, not material, but in the order of efficient cause, which must be most perfect. For just as matter, as such, is merely potential, an agent, as such, is in the state of actuality. Hence, the first active principle must needs be most actual, and therefore most perfect; for a thing is perfect in proportion to its state of actuality, because we call that perfect which lacks nothing of the mode of its perfection.
\end{quote}

\subsubsection*{Claim}

Aquinas asserts that because God is the first efficient cause rather than a material principle, and an active agent is actual as such, God must be maximally actual and therefore maximally perfect, since perfection tracks actuality.

\emph{Plain-language reading.} Whereas raw matter is incomplete capacity, an active cause operates insofar as it is actual. The first active cause must therefore be maximally actual, and because we call something perfect insofar as it lacks none of the actuality proper to its kind, it must be maximally perfect.

\subsubsection*{Meaning}

\emph{Perfect} here does not mean ``morally virtuous'' or ``pleasing to human taste,'' nor ``completely made'' (which Aquinas drops in Reply 1, extending the word to whatever lacks nothing in actuality even if unmade). \emph{Actuality} means realized ontological completeness, opposed to unfinished potentiality.

\subsubsection*{Warrant}

Matthew 5:48 in the \emph{sed contra} stays labelled as \textbf{revelation}, outside the philosophical argument of the body. Inside the body, ``God is the first principle\dots\ in the order of efficient cause'' is a \textbf{prior result} from \emph{ST} I, q.\,2, a.\,3. That matter is potential while an agent is actual pairs Aristotelian \textbf{definition} with a \textbf{metaphysical principle}; that the first active principle is maximally actual enters as a \textbf{logical consequence} leaning on q.\,3, a.\,1's claim that the first being has no potentiality. The final clause---``a thing is perfect in proportion to its state of actuality''---reads like a \textbf{definition}, yet bridges an ontological term (actuality) to an evaluative term (perfection).

\subsubsection*{Burden}

Granting the prior results from Questions 2 and 3, the burden-bearing premise is that closing bridge: \emph{that ontological actuality is the objective measure of perfection}. An intelligent opponent can accept what ``efficient cause'' and ``actuality'' mean and grant a maximally actual first cause while denying that sheer ontological actuality amounts to perfection in any evaluative sense.

\subsubsection*{Consequence}

Granting the prior results and that bridge establishes that the first active principle lacks none of the ontological actuality proper to its mode of being.

\emph{Not yet established.} This article has not established that all created perfections pre-exist in God (argued in q.\,4, a.\,2), nor that God is morally good, loving, merciful, or just (argued separately in qq.\,5--6 and qq.\,19--21).

\subsubsection*{Denial}

Here \textbf{non-assent} and \textbf{rival framework} separate cleanly from \textbf{contradiction}. If a careless defender declares that denying pure act is perfect is ``self-contradictory because perfection just means fullness of act,'' he mistakes refusal of a stipulated vocabulary for a broken law of logic.

In its simplest form, declining to equate perfection with actuality is \textbf{non-assent}. A reader may say: ``I see how your Aristotelian vocabulary uses \emph{perfectum} to track fullness of act; I simply decline that usage, because for me `perfect' is an evaluative word for excellence or fitness for purpose, not degree of actuality.'' He incurs no contradiction and contradicts no experience. Yet he \textbf{gives up} something real: he forfeits an objective metaphysical measure tying being and goodness together, leaving himself to account elsewhere for what makes anything genuinely better than anything else.

A thinker who constructs that alternative account enters a \textbf{rival framework}---a developed \emph{teleological or evaluative account} on which perfection is always relative to functional excellence, flourishing, or moral standards rather than a single ontological scale of actuality. On that rival view, a completely actual yet impersonal cosmic force would lack nothing in actuality without being ``perfect'' in any evaluative sense worthy of reverence. The disagreement is now a substantive clash over whether goodness is rooted in being itself or arises only with life, mind, and purpose.

\subsection{Exercises}

\subsubsection*{Exercise 5.1}
Read the following short argument and answer the questions beneath it.

\begin{quote}
Whenever dry salt is stirred into room-temperature water, it dissolves reliably. Bare geometric positions of atoms at one instant do not logically entail their positions at the next instant. Therefore, salt and water possess real, irreducible causal powers---intrinsic active and passive capacities beyond their static arrangement---that produce dissolution on contact.
\end{quote}

Identify the smallest denial that blocks the conclusion while granting as much as possible, classify it accurately among the five kinds of denial, and explain why two other kinds do not fit.

\subsubsection*{Exercise 5.2}
Read the following passage, written in the voice of a scholastic manual, and test its closing verdict.

\begin{quote}
Every changing bodily substance is necessarily composed of a real underlying substance and distinct accidental forms. For by ``substance'' we mean the enduring subject that stays one and the same while gaining or losing qualities, and by ``accidental form'' we mean a real determination that comes and goes in it, as an oak turns from green in summer to russet in autumn while remaining the same tree. Anyone who denies this real distinction therefore contradicts himself: he admits that the oak changes colour while denying the very principle that makes change in a single subject possible.
\end{quote}

Does denying the burden-bearing premise here incur a \textbf{contradiction}? Classify the denial accurately, explain in full why the denier incurs no contradiction, and state clearly what he gives up.

\subsubsection*{Exercise 5.3}
Read this unannotated paragraph from a modern critic of Aquinas's First Way.

\begin{quote}
Aquinas's First Way collapses on the first page of modern physics. Aquinas assumes that whatever is in motion must be continually pushed along by an external mover right now, as an ox drags a cart through mud and the cart stops the instant the ox halts. Newton's First Law exploded that medieval assumption three centuries ago: a body in uniform straight-line motion continues moving indefinitely with no external force acting on it. Since inertial coasting refutes Aquinas's central premise by direct counterexample, the argument from motion is dead.
\end{quote}

Diagnose the exact point where this counterexample fails to touch Aquinas's claim in \emph{ST} I, q.\,2, a.\,3, corp., and identify the sharper inertial objection that \emph{does} still press Aquinas.

\subsubsection*{Exercise 5.4}
Read this unannotated defective argument and its reply to an imaginary critic.

\begin{quote}
Every natural agent acts for an intrinsic final end, because to ``act'' simply means to strive purposefully toward a built-in natural goal. If a critic objects that rivers flow downhill and enzymes catalyze reactions purely through mechanical interaction without intrinsic goals, the answer is plain: to deny that a natural agent aims at an end is to deny that it acts at all, contradicting our everyday observation that rivers flow and enzymes work.
\end{quote}

Diagnose both the definitional smuggling and the pseudo-reply. Then repair the argument so that it makes a candid case for natural finality \emph{without} smuggling the conclusion into a definition, prove via the hostile-reader test that the repaired premise leaves a genuine substantive gap open, and show one tempting repair that cheats.

\subsubsection*{Exercise 5.5}
Read this unannotated excerpt from St.\ Thomas Aquinas, \emph{Summa Theologiae} I, q.\,3, a.\,2, corp.\ (``Whether God Is Composed of Matter and Form?'', 1920 English Dominican translation). Nothing has been marked or supplied.

\begin{quote}
It is impossible that matter should exist in God. First, because matter is in potentiality. But we have shown (Q.\,2, A.\,3) that God is pure act, without any potentiality. Hence it is impossible that God should be composed of matter and form. Secondly, because everything composed of matter and form owes its perfection and goodness to its form; therefore its goodness is participated, inasmuch as matter participates the form. Now the first good and the best---viz.\ God---is not a participated good, because the essential good is prior to the participated good. Hence it is impossible that God should be composed of matter and form. Thirdly, because every agent acts by its form; hence the manner in which it has its form is the manner in which it is an agent. Therefore whatever is primarily and essentially an agent must be primarily and essentially form. Now God is the first agent, since He is the first efficient cause. He is therefore of His essence a form; and not composed of matter and form.
\end{quote}

Identify the claim, locate the burden-bearing premises (checking any promissory references), note what remains not yet established, classify denial of the hylomorphic premises among the five kinds, and state the strongest rival.

\subsection{Solutions}

\subsubsection*{Solution to Exercise 5.1}
The \textbf{smallest denial} grants both stated premises (salt reliably dissolves, and static snapshots do not logically entail future positions) and rejects the hidden burden-bearing premise that \emph{lawful regularity between successive states requires real, irreducible causal powers residing inside the substances}. That denial is a \textbf{rival framework} paired with an \textbf{explanatory cost} dispute: a Humean regularity ontology or a governing-laws account explains dissolution without intrinsic Aristotelian powers. It is not \textbf{contradiction} (no contradiction follows from lawful sequences without inner powers) and not \textbf{conflict with experience} (we observe dissolution, not an invisible power). Classifying it solely as \textbf{explanatory cost} catches the Humean's debt to explain regularity, yet misses that serious Humeans and governing-law realists defend complete competing frameworks while pressing their own explanatory cost back on irreducible powers.

\subsubsection*{Solution to Exercise 5.2}
No: denial incurs zero \textbf{contradiction}. The passage stipulates that ``substance'' means an underlying substratum distinct from accidental forms, then accuses anyone refusing that split of contradicting himself when calling the autumn tree the same oak. Construct an explicit denier: ``I agree that this oak is green in July and russet in October and that ordinary speech calls it the same tree; I simply decline your vocabulary of an enduring substantial substratum distinct from accidental forms.'' He never affirms and denies the same claim in the same respect, because he never accepted the stipulated equation of enduring sameness with substance-accident composition.

In its minimal form, declining that idiom without building a competing ontology is \textbf{non-assent} (if the denier goes further and constructs a four-dimensionalist temporal-parts or trope-bundle ontology, it becomes a \textbf{rival framework}---both verdicts are defensible depending on how far the denier goes). Crucially, the non-assenter \textbf{gives up} a straightforward ontological anchor for strict numerical identity across change (why sapling and mature oak are literally one enduring individual rather than connected temporal stages) and a principled distinction between what a thing essentially is and what merely happens to it.

\subsubsection*{Solution to Exercise 5.3}
The critic fails counterexample discipline because Aquinas neither mentions an ox pulling a cart nor defines motion as uniform travel requiring a continuous push. In \emph{ST} I, q.\,2, a.\,3, corp., Aquinas explicitly defines \emph{motus} as ``the reduction of something from potentiality to actuality,'' illustrated by fire heating wood. In Newtonian physics, uniform straight-line velocity is a conserved state rather than an acquisition of a new quality or acceleration; the critic's counterexample attacks a caricature of medieval impetus physics while leaving heating, chemical alteration, and acceleration untouched.

A \emph{sharper} inertial objection genuinely presses Aquinas, however, if a defender insists that a coasting body continuously actualizes fresh spatial potencies at every new point along its path. On that reading, Newton's First Law directly challenges ``whatever is in motion is put in motion by another,'' forcing either a clash of \textbf{rival frameworks} over whether spatial location is an accidental form continuously actualized by a concurrent cause, or a candid Thomistic concession that uniform inertial velocity is a state rather than \emph{motus}, narrowing the First Way to genuine potency-to-act changes.

\subsubsection*{Solution to Exercise 5.4}
The opening sentence smuggles the conclusion into a loaded \textbf{definition} of ``act'' (striving purposefully toward an intrinsic goal), turning the argument into a verbal circle. The second sentence adds a \textbf{pseudo-reply}, translating the critic's mechanistic objection into that loaded definition and falsely accusing him of denying that rivers flow.

An honest repair separates definition from substantive principle: ``To `act,' neutrally defined, means to produce a determinate physical effect. Experience shows that non-intelligent natural things---acorns growing into oaks, enzymes catalyzing specific reactions---regularly produce determinate ranges of effects rather than random outcomes. Yet such reliable directedness cannot be accounted for merely by antecedent events or bare regularities; it requires real intrinsic tendencies in natural agents oriented toward their characteristic effects. Therefore, non-intelligent natural agents act for intrinsic ends.''

Run the \textbf{hostile-reader test}: an opponent can accept the neutral definition of ``act'' and every observed regularity about acorns and enzymes, yet coherently deny the substantive finality principle (adopting a Humean or purely structural-law \textbf{rival framework}). Because the gap stays open, the repaired premise is genuinely substantive. A tempting repair that cheats would redefine ``regular physical effect'' as ``an effect teleologically aimed at by a thing's essence,'' rebuilding the tautology under new words.

\subsubsection*{Solution to Exercise 5.5}
Aquinas's \textbf{claim} is that God cannot be composed of matter and form, because God is pure act, essential good, and primary agent, and therefore essentially form. Three \textbf{prior results} enter here as promissory notes: pure act without potentiality, first good, and first efficient cause. Inspect the first citation carefully: Aquinas points to ``(Q.\,2, A.\,3)'' for ``pure act, without any potentiality,'' yet q.\,2, a.\,3 argued only to an unmoved mover---it was q.\,3, a.\,1 that actually argued God has no potentiality!

Even granting those prior results, three hylomorphic \textbf{metaphysical principles} bear the fresh burden inside this paragraph: (a) matter as such is potentiality; (b) a composite's goodness is merely participated through matter receiving form; and (c) every agent acts through its form, so a primary agent must be pure form essentially. Even if granted, the paragraph has \textbf{not yet established} that God is identical with His existence (q.\,3, a.\,4), possesses intellect, consciousness, or free will (later questions), or created matter from nothing (q.\,44, a.\,2).

Denying these hylomorphic premises incurs neither \textbf{contradiction} nor bare \textbf{conflict with experience}. Declining matter-and-form talk altogether is \textbf{non-assent}; adopting a developed alternative is a \textbf{rival framework}. At its strongest, that rival rejects equating matter with inert passive potentiality, arguing from dynamic physical-field ontologies (or non-hylomorphic philosophical theism) that physical matter-energy is intrinsically active rather than passive receptivity, that causation is not form impressing passive matter, and that proving a first cause immaterial must rest directly on the contingency and finitude of physical reality rather than on Aristotelian hylomorphism. (A reader who instead targets the inherited promissory notes from Question 2 makes an equally defensible move, though isolating the hylomorphic premises exposes this article's own burden.)
